%% preamble.tex --- パッケージ読み込みと共通マクロ

%% ---------- フォント・基本 ----------
\usepackage{luatexja-fontspec}
\usepackage{amsmath,amssymb}
\usepackage{graphicx}
\usepackage{flafter}  % 図は本文で書かれた位置より後ろにだけ置く（章のタイトルより上に来ないように）
\usepackage{tikz}  % 図版（images/ 以下の .tex）とプレースホルダ用
\usetikzlibrary{positioning,arrows.meta,fit,calc,backgrounds,automata,shapes.geometric}
% 図で共通して使うスタイル
\tikzset{
  figbox/.style={draw, rounded corners=2pt, align=center, font=\small,
    minimum height=8mm, inner sep=3pt, fill=white},
  figpart/.style={figbox, fill=blue!6},
  figarrow/.style={-{Stealth[length=2mm]}, thick},
  figlabel/.style={font=\footnotesize, align=center},
  fignum/.style={circle, draw, fill=white, inner sep=0.8pt, font=\scriptsize\bfseries},
}
\usepackage{xcolor}
\usepackage{booktabs}
\usepackage{url}
\usepackage[hidelinks]{hyperref}
\graphicspath{{images/}}

%% ---------- コードリスト ----------
\usepackage{listings}
\lstset{
  basicstyle=\ttfamily\small,
  columns=fullflexible,
  keepspaces=true,
  breaklines=true,
  frame=single,
  framerule=0.4pt,
  rulecolor=\color{black!40},
  backgroundcolor=\color{black!3},
  numbers=left,
  numberstyle=\tiny\color{black!50},
  % 行番号を枠の中に入れる（枠の左側に行番号の分の余白を取る）
  numbersep=0.8em,
  xleftmargin=2.2em,
  framexleftmargin=1.8em,
  keywordstyle=\color{blue!70!black}\bfseries,
  commentstyle=\color{green!40!black},
  stringstyle=\color{red!60!black},
  showstringspaces=false,
  tabsize=4,
}
\lstdefinestyle{python}{language=Python}
\lstdefinestyle{bash}{language=bash}
\lstdefinestyle{cpp}{language=C++}
\lstdefinestyle{xml}{language=XML}

%% ---------- 囲み ----------
\usepackage[most]{tcolorbox}

% ターミナル操作例:  \begin{terminal}[タイトル] $ ls -l \end{terminal}
% Ubuntu（GNOME）のターミナル風：中央にタイトル、右上にウィンドウのボタン、明るい本文。
% `$` から行末まで（入力するコマンド）を太字にする。
\definecolor{ubuntuorange}{HTML}{E95420}
\definecolor{ubuntutitlebar}{HTML}{3A3A3A}
\newcommand{\terminalbuttons}{%
  \begin{tikzpicture}[baseline=-0.6ex, line width=0.5pt]
    \fill[white!55!ubuntutitlebar] (0,0) circle (1.5mm);
    \draw[white] (-0.7mm,0) -- (0.7mm,0);
    \fill[white!55!ubuntutitlebar] (4.2mm,0) circle (1.5mm);
    \draw[white] (3.5mm,-0.7mm) rectangle (4.9mm,0.7mm);
    \fill[ubuntuorange] (8.4mm,0) circle (1.5mm);
    \draw[white] (7.8mm,-0.6mm) -- (9.0mm,0.6mm) (7.8mm,0.6mm) -- (9.0mm,-0.6mm);
  \end{tikzpicture}}
\newtcblisting{terminal}[1][]{
  listing only, enhanced,
  colback=black!3, colframe=ubuntutitlebar, arc=2pt, boxrule=0.6pt,
  left=5pt, right=4pt, top=2pt, bottom=2pt,
  listing options={basicstyle=\ttfamily\small,
    columns=fullflexible, keepspaces=true, breaklines=true,
    numbers=none, frame=none, backgroundcolor={}, aboveskip=0pt, belowskip=0pt,
    xleftmargin=0pt, framexleftmargin=0pt,  % ソースコード用の行番号の余白を打ち消す
    morecomment={[l][\bfseries]{\$}}},
  title={\makebox[\linewidth][c]{\strut#1}},
  fonttitle=\sffamily\small\bfseries, colbacktitle=ubuntutitlebar, coltitle=white,
  toptitle=0.8pt, bottomtitle=0.8pt,
  overlay={\node[anchor=east, inner sep=0pt] at ([xshift=-2.5mm]title.east) {\terminalbuttons};},
}

% 表示例のターミナル（読者が入力しない例）：紺色のタイトルバー・点線の枠・タイトルの先頭に「例：」
%   \begin{termexample}[タイトル] ... \end{termexample}
% 存在しないファイルや機器を使う例、実行すると困る例（kill など）、説明のための書式の例に使う。
\definecolor{examplebar}{HTML}{2B4C7E}
\newtcblisting{termexample}[1][]{
  listing only, enhanced,
  colback=examplebar!4, colframe=examplebar, arc=2pt, boxrule=0pt,
  borderline={0.8pt}{0pt}{examplebar, dashed},
  left=5pt, right=4pt, top=2pt, bottom=2pt,
  listing options={basicstyle=\ttfamily\small,
    columns=fullflexible, keepspaces=true, breaklines=true,
    numbers=none, frame=none, backgroundcolor={}, aboveskip=0pt, belowskip=0pt,
    xleftmargin=0pt, framexleftmargin=0pt,
    morecomment={[l][\bfseries]{\$}}},
  title={\makebox[\linewidth][c]{\strut 例：#1}},
  fonttitle=\sffamily\small\bfseries, colbacktitle=examplebar, coltitle=white,
  toptitle=0.8pt, bottomtitle=0.8pt,
}

% ページをまたいだ囲みに「続く」を表示する（breakable な囲みで使う）
\usepackage{needspace}
\tikzset{contmark/.style={font=\scriptsize\sffamily, text=black!55, inner sep=1.5pt}}
\tcbset{
  contmarks/.style={
    pad at break*=3.5mm,
    overlay first={\node[contmark, anchor=south east] at (frame.south east) {（次のページに続く）};},
    overlay middle={\node[contmark, anchor=south east] at (frame.south east) {（次のページに続く）};
      \node[contmark, anchor=north east] at (frame.north east) {（前のページから続く）};},
    overlay last={\node[contmark, anchor=north east] at (frame.north east) {（前のページから続く）};},
  },
}
% ページの残りが少ないときは、囲みごと次のページに送る（数行だけ前のページに残らないように）
\newcommand{\keeplines}{\par\needspace{6\baselineskip}}

% ソースコードの枠（キャプション付き）：ページをまたぐと「続く」を表示する
%   \begin{codelisting}[listings のオプション]{キャプション}{ラベル} ... \end{codelisting}
%   \codefile[listings のオプション]{キャプション}{ラベル}{ファイル名}
% キャプションのないコード（囲みの中の短い例など）は、これまでどおり lstlisting を使う。
\tcbset{
  codebox/.style={
    enhanced, breakable, contmarks,
    colback=black!3, colframe=black!40, boxrule=0.4pt, arc=1pt,
    left=2pt, right=4pt, top=2pt, bottom=2pt,
    colbacktitle=black!10, coltitle=black, fonttitle=\small, halign title=flush center,
    toptitle=0.5pt, bottomtitle=0.5pt, titlerule=0.4pt,
  },
  codelisting options/.style={listing options={
    frame=none, backgroundcolor={}, aboveskip=0pt, belowskip=0pt,
    xleftmargin=1.8em, framexleftmargin=0pt, numbersep=0.8em, #1}},
}
\newtcblisting[auto counter, number within=chapter,
  number freestyle={\noexpand\chapprefix\noexpand\thechapter.\noexpand\arabic{\tcbcounter}}]
  {codelisting}[3][]{
  codebox, listing only, codelisting options={#1},
  title={Listing \thetcbcounter\quad #2}, label={#3},
}
\AddToHook{env/codelisting/before}{\keeplines}
\newtcbinputlisting[use counter from=codelisting]{\codefileinner}[4][]{
  codebox, listing only, codelisting options={#1}, listing file={#4},
  title={Listing \thetcbcounter\quad #2}, label={#3},
}
\newcommand{\codefile}{\keeplines\codefileinner}

% この章で学ぶこと
\newtcolorbox{goalbox}{
  colback=blue!4, colframe=blue!50!black, arc=2pt,
  title={この章で学ぶこと}, fonttitle=\sffamily\bfseries,
}

% ポイント・注意・コラム
\newtcolorbox{point}[1][ポイント]{
  colback=green!4, colframe=green!40!black, arc=2pt,
  title={#1}, fonttitle=\sffamily\bfseries,
}
\newtcolorbox{caution}[1][注意]{
  colback=red!4, colframe=red!60!black, arc=2pt,
  title={#1}, fonttitle=\sffamily\bfseries,
}
\newtcolorbox{column}[1]{
  colback=orange!4, colframe=orange!60!black, arc=2pt,
  title={コラム：#1}, fonttitle=\sffamily\bfseries,
}

% TIPS：本筋とは別の、すぐに役立つ実践的なアドバイス（長くてもページをまたげる）
%   \begin{tips}{タイトル} ... \end{tips}
\newtcolorbox{tips}[1]{
  enhanced, breakable, contmarks,
  colback=teal!4, colframe=teal!60!black, arc=2pt,
  title={TIPS：#1}, fonttitle=\sffamily\bfseries,
}
\AddToHook{env/tips/before}{\keeplines}

% 本文中の表：前後に余白を空け、表の上にキャプション（表 5.1 など）を付ける
%   \begin{tablebox}[キャプション][tab:ラベル] \begin{tabular}{ll} ... \end{tabular} \end{tablebox}
%   キャプションを省略すると、番号なしの表になる
\usepackage{capt-of}
\NewDocumentEnvironment{tablebox}{o o}{%
  \par\addvspace{\bigskipamount}%
  \noindent\begin{minipage}{\linewidth}% キャプションと表を同じページに置く
  \centering
  \IfValueT{#1}{\captionof{table}{#1}\IfValueT{#2}{\label{#2}}}%
}{%
  \end{minipage}%
  \par\addvspace{\bigskipamount}%
}

% 練習問題：\begin{exercise}[label={ex:...}] \begin{enumerate} \item ... \end{enumerate} \end{exercise}
% 章ごとに「練習問題 18.1」のように番号が振られ、中の問題は (1), (2), ... になる。ページをまたげる。
\newtcolorbox[auto counter, number within=chapter]{exercise}[1][]{
  enhanced, breakable, contmarks,
  colback=violet!3, colframe=violet!55!black, arc=2pt,
  title={練習問題 \thetcbcounter}, fonttitle=\sffamily\bfseries,
  before upper={\renewcommand{\labelenumi}{(\arabic{enumi})\hspace{0.5em}}\setlength{\leftmargini}{2.8em}},
  #1
}
\AddToHook{env/exercise/before}{\keeplines}

%% ---------- 共通マクロ ----------
\newcommand{\cmd}[1]{\texttt{#1}}          % インラインのコマンド・ファイル名
\newcommand{\key}[1]{\fbox{\small\texttt{#1}}} % キー入力 例: \key{Ctrl}+\key{C}
\newcommand{\todo}[1]{{\color{red}[TODO: #1]}}

%% ---------- 補章 ----------
% 補章：本編の章の内容を、別の書き方などで補う章（付録は早見表・用語集などの資料）
% main.tex で \hoshou と書くと、その後の \chapter が「補章1」「補章2」……になる。
% 節・図・表・リストの番号は「補1.1」のようにして、本編の章（1.1 など）と区別する。
% 本文から参照するときは「補章\ref{sup:...}」と書く。
\newcommand{\chapprefix}{}
\renewcommand{\thesection}{\chapprefix\thechapter.\arabic{section}}
\renewcommand{\thefigure}{\chapprefix\thechapter.\arabic{figure}}
\renewcommand{\thetable}{\chapprefix\thechapter.\arabic{table}}
\AtBeginDocument{\renewcommand{\thelstlisting}{\chapprefix\thechapter.\arabic{lstlisting}}}
\newcommand{\hoshou}{%
  \clearpage
  \setcounter{chapter}{0}%
  \renewcommand{\thechapter}{\arabic{chapter}}%
  \renewcommand{\chapprefix}{補}%
  \ModifyHeading{chapter}{label_format={補章\thechapter}}%
  \def\chaptermark##1{\markboth{補章\thechapter\quad ##1}{}}%
}
% \appendix で、番号の付け方を本編と同じ形に戻す
\jlreqsetup{appendix_precode={\renewcommand{\chapprefix}{}}}
